% Section 1 -- Introduction

A bi-level recursive self-improvement (RSI) system is one in which an \emph{outer-loop} agent $\Omega$ rewrites the code of an \emph{inner-loop} agent $\Lambda$, with retention decisions driven by a private evaluation signal $\mu_{\text{priv}}$ that is invisible to $\Lambda$ itself. The archetype was operationalized publicly in 2026 by Weco AI's AIDE\textsuperscript{2}, in which an eight-day, 100-iteration outer loop (driven by a hand-tuned autonomous research agent using a frontier LLM) rewrote a simplified AIDE agent (driven by a smaller LLM) seven times, with rewritten agents retained only when their held-out private score exceeded the previous best. Roughly nine in ten proposed rewrites were rejected.

This pattern is not unique to AIDE\textsuperscript{2}. Any system that couples an \emph{artifact generator} (inner loop) with a \emph{generator-of-the-generator} (outer loop) under a selection pressure grounded in hidden evaluation exhibits the structure: RL-based prompt search with private evaluation harnesses, evolutionary architecture search with held-out benchmarks, automated research systems targeting peer-review-grounded feedback. We argue in this paper that the literature lacks a unified formal vocabulary for reasoning about such systems, and that this gap is consequential for both scientific understanding and engineering deployment.

\subsection{The Audit Channel Gap}

In every bi-level RSI system of which we are aware, including AIDE\textsuperscript{2}, the \textbf{audit channel} between the inner loop's optimization signal and the outer loop's selection signal is treated as implementation plumbing rather than as a load-bearing architectural component. This is a category error with predictable consequences: AIDE\textsuperscript{2}'s published LLM context lacks an explicit isolation mechanism ($\mu_{\text{priv}}$ can in principle leak through prompt-injection vectors), the outer loop's mutation of its own eval pipeline (the ``eval-repair monkey patch'' of iteration $\sim$70) is not auditable, and the audit trail of L3 statistical-layer changes is append-only in name only. We argue this is a \emph{category-level} gap in the bi-level RSI literature, not a bug in any specific system.

We define the audit channel $\alpha$ as a first-class object --- a sequentially-ordered, content-hashed log of system state transitions --- and show in \S3.3 how this definition captures three properties current systems implement only by accident. \S3 then develops the formal framework, \S4 instantiates it on AIDE\textsuperscript{2}, and \S5 supplies architectural recommendations.

\subsection{Contributions}

This paper makes four contributions:

\begin{enumerate}
    \item \textbf{A formal framework} (\S3) that elevates the audit channel to a first-class component of bi-level RSI systems and instantiates it as a Policy-Evaluation Pair Automaton (PEPA) with three explicit audit-channel properties.
    \item \textbf{Three falsifiable predictions} (\S3.4) that distinguish our framework from informal accounts of bi-level RSI: (P1) a closed-form effect-size bound for ignition tests under measurement noise, (P2) a proof-of-concept that $\mu_{\text{priv}}$ must be ungameable for Level-2 RSI, and (P3) an alternative explanation for observed anti-reward-hacking emergence.
    \item \textbf{An empirical instantiation} (\S4) on AIDE\textsuperscript{2} with a reference implementation of P2 (open-sourced as part of the paper's reproducibility package). On synthetic data calibrated to AIDE\textsuperscript{2}'s reported statistics, our P2 meta-predictor achieves canonical single-split $\text{AUC} = 0.9970$ (10-fold CV mean $0.9957\pm0.0027$), falsifying P2 in a controlled setting; this synthetic test is \emph{structurally circular} and is demoted to a directional demonstration with a specified non-circular audit-gate fix (\S4.4); real AIDE85 lineage is awaited for empirical confirmation.
    \item \textbf{Architectural recommendations} (\S5) for closing the audit-channel gap: cryptographic commitment to $\mu_{\text{priv}}$, multi-PDP voting on rewrite acceptance, $\Omega$ access-audit logging, and TEE-mediated channel separation. Each recommendation is framed as a minimal viable extension that an existing system could adopt incrementally.
\end{enumerate}

We position the paper as a \emph{methods paper} in TMLR's submission taxonomy: a formal framework plus a single concrete instantiation plus open empirical questions (see \S7.4 for detailed positioning). The framework is descriptive rather than constructive: it offers vocabulary, falsifiable predictions, and architectural recommendations, but does not itself implement an improved bi-level RSI system. We see this as a strength; the gap between descriptive frameworks and deployed systems has historically been productive in access-control research and we expect the same here.

We explicitly do not claim to have reproduced AIDE\textsuperscript{2} or to have access to its unpublished components. All AIDE\textsuperscript{2}-specific analysis is based on the public blog post (Weco AI, 2026) and on AIDE0's open-sourced code (Weco AI, 2025, arXiv:2502.13138). Where our claims require lineage data we have not been granted, we say so.

\subsubsection{Quick Reference for Reviewers}

For reviewers familiar with access-control engineering but new to bi-level RSI, the following summary may be useful. \textbf{Key concept}: the \emph{audit channel} $\alpha$ between inner and outer loops is the load-bearing object; we argue that any bi-level RSI without a properly specified $\alpha$ cannot reach Level-2 regardless of inner-loop sophistication. \textbf{Three falsifiable predictions}: P1 (Appendix~C) gives a closed-form noise bound for the ignition test, with proof under i.i.d.\ and AR(1) evaluation models; P2 (Appendix~D) shows that predictable $\mu_{\text{priv}}$ is sufficient to block Level-2 RSI, via three lemmas establishing gap monotonicity; P3 (Appendix~E) defines audit-channel capacity $C(\alpha)$ and derives the relation $r(\alpha) = 1 - (C(\alpha) - C_0)/(C_{\max} - C_0)$ under TEE-isolated audit. \textbf{Empirical result}: the P2 meta-predictor achieves canonical single-split $\text{AUC} = 0.9970$ on synthetic data calibrated to AIDE\textsuperscript{2}'s reported statistics (\S4.4), with 5-model 10-fold robustness confirming the result is not a linear-model artifact --- \emph{subject to the non-circularity caveat of \S4.4}. \textbf{Four architectural recommendations} (R1--R4) provide concrete mitigations, with R1 (cryptographic commitment to $\mu_{\text{priv}}$) as the highest-signal-to-effort first step. \textbf{Companion code} at the anonymized companion repository reproduces all reported numbers in $<1$ minute on commodity hardware.
